% !PRINTED 387
\weberpage{387}

%<<CONT
the interchange of $[1\,4]$ with $[1\,5]$ in agreement with formula (10). Just as easily one recognises that, e.g., $[1\,2]$ is not altered by (12).

\begin{center}
\S.\,103.

\gesperrt{Simplicity of the group of the double tangent problem.}
\end{center}

The representation of the group $P$ which we developed in the preceding section now furnishes us a quite simple proof that this group has \gesperrt{no normal divisor}, and hence, in the manner of expression we used earlier, is \gesperrt{simple}; from which it then follows that the group cannot be lowered by the adjunction of irrationalities with smaller group, and thus, for example, not by cyclic equations.

We represent $P$ in the form (4), \S.\,102:
\weq{1}{P = S + \Sigma\, S\, \Pi_{\alpha_1, \alpha_2, \alpha_3, \alpha_4},}
in which $S$ is the whole group of all permutations of eight figures. The group $S$ has a normal divisor $S'$ of index 2, namely the alternating group of the eight figures, which for its part is simple, and every normal divisor of $S$ that does not consist of the identical permutation alone must contain the whole group $S'$. We set
\weq{2}{S = S' + S'',}
in which $S'' = S'\sigma$, where $\sigma$ denotes any permutation of the second kind, e.g. a transposition.

We now assume that $Q$ is a normal divisor of $P$ which does not consist of the identical substitution alone.

The greatest common divisor of $Q$ and $S$ is then a normal divisor of $S$, and must therefore, if it is not the identical group, contain the group $S'$. From this it follows:

1.\ \gesperrt{If $Q$ contains a non-identical permutation from $S$, then $Q$ contains the whole group $S'$.}

We then prove that $Q$, if it contains the group $S'$, must be identical with $P$.

For if $Q$ contains the whole group $S'$, then it contains
%CONT>>
